\subsection{A critical estimate on an entire marked shell}
\label{sec:critical-annulus}

The low-curvature shell between separated quantiles has small
\emph{physical} full curvature. The following criterion uses that
smallness on one actual coordinate domain. Write
\(A(c,C)=\{c<|x|<C\}\subset\mathbb R^4\).

\begin{proposition}[Harmonic annulus at fixed distortion]
\label{prop:critical-annulus}
There is a fixed \(\varepsilon_*>0\) with the following property.
On the separated subsequence above, suppose there are oriented smooth
embeddings of \(A(1/8,8)\) into the intervening physical shells,
with the inner charged ball and the entire median exterior on
opposite sides of the middle sphere and both disjoint from the entire closed trim \(\{1/4\le|x|\le4\}\). Suppose that in these coordinates
a constant rescaling \(m_i\) of the physical metric satisfies
\[
 \Delta_{m_i}x^a=0,\qquad
 \|m_i-\delta_{\rm Euc}\|_\infty\le\varepsilon_*,
 \qquad \int_{A(1/8,8)}|\Rm_{m_i}|^2\,dV_{m_i}\longrightarrow0 .
\]
This configuration is impossible on the fixed nonstandard type at
its positive-class infimum. The distortion bound is fixed before
taking the source limit. The same conclusion holds for approximately harmonic coordinates with \(\|D(\Delta_{m_i}x)\|_{L^2}\to0\) on this fixed annulus.
\end{proposition}

\begin{proof}
We give the critical estimate in the coordinates where the metric is
known. For a smooth \(m=\delta_{\rm Euc}+v\) write
\(\eta=\|v\|_\infty\) and \(T^a=\Delta_m x^a\).
For a fixed compactly supported cutoff \(0\le\chi\le1\) put
\[
 X=\|\chi^2D^2v\|_2,\quad G=\|\chi Dv\|_4,\quad
 L_\chi=\|D\chi\|_4^2+\|D^2(\chi^2)\|_2 .
\]
Norms of components use \(dx\). Integrating
\(\chi^4|Dv|^4\) by parts against \(v\) gives
\(G^4\le C\eta(G^2 X+\|D\chi\|_4G^3)\), hence
\(G^2\le C\eta X+C\eta^2 L_\chi\).
The coordinate Ricci equation has principal part
\[
 m^{ab}\partial_{ab}v_{jk}
 =-2\Ric_{jk}-m_{ja}\partial_kT^a-m_{ka}\partial_jT^a
     +O(|Dv|^2).
\]
Apply \(\|D^2f\|_2=\|\Delta_{\rm Euc}f\|_2\) to
\(f=\chi^2v\) extended by zero. The commutators cost
\(C\|D\chi\|_4G+C\eta L_\chi\), and freezing the principal
coefficient costs \(C\eta X\). Absorbing these terms using
the bound for \(G^2\) gives
\begin{equation}\label{eq:critical-chart}
 X\le C B,\qquad G^2\le C\eta B,\qquad
 B=\|\chi^2\Ric\|_2+\|\chi^2 DT\|_2+\eta L_\chi .
\end{equation}
In particular if \(\eta_i\to0\) and \(B_i\) stays bounded, this
gives the strong first-order estimate and bounded full curvature on
each interior splice. The differentiated tension belongs to the
coordinate system used in the interpolation.

For the fixed-distortion assertion choose cutoffs on nested annuli
in \(A(1/8,8)\) which equal one on a neighborhood of
\(\overline{A(1/4,4)}\). In harmonic coordinates \(T=0\); in the approximate case the differentiated tension contributes an additional \(o(1)\).
Equation~\eqref{eq:critical-chart} gives bounded \(W^{2,2}\) and
\(\|Dm_i\|_4\le C\varepsilon_*+o_i(1)\) there. Pass to a
weak \(W^{2,2}\) limit \(h\), strongly in each \(W^{1,s}\),
\(s<4\), and almost everywhere. The matrices and inverses stay
uniformly within \(C\varepsilon_*\) of the Euclidean ones.
In the Ricci and harmonic equations the bounded coefficients
converge almost everywhere. They may be transferred to a fixed
test in \(L^2\) in the Hessian term. Quadratic gradient terms
converge in distributions by strong \(W^{1,2}\) and the uniform
ellipticity. Thus the limit satisfies the reduced metric equation with vanishing right-hand source Ricci and differentiated tension. In the approximate case a limiting constant coordinate tension is allowed; its derivative vanishes throughout the argument. Its full curvature also vanishes in distributions,
by the all-lowered formula and smallness of the incoming integral.

Here is the interior regularity at the critical exponent. Write
the equation for \(h\) as
\(h^{ab}\partial_{ab}h+C(h,Dh)Dh=0\), where \(C\) denotes a
matrix of first-order coefficients bounded by \(C_0|Dh|\).
On a fixed ball compactly inside the annulus, this system with
zero Dirichlet data is a small operator perturbation of the flat
Laplacian from \(W^{2,s}\cap W^{1,s}_0\) to \(L^s\), for
\(s=2,3\). Indeed the principal error is \(O(\varepsilon_*)\);
the first-order error is bounded by
\(C\|Dh\|_4\|Du\|_{4s/(4-s)}\).
Choose \(\varepsilon_*\) for these two estimates and the fixed
cutoffs. Applying the operator to a compactly supported cutoff
times \(h\) gives a forcing in \(L^3\), since \(Dh\in L^4\).
Solve at exponent three and use uniqueness at two to identify
the solution. Thus \(h\in W^{2,3}_{\rm loc}\); its reduced metric equation first gives \(W^{2,6}\) by \(Dh\in L^{12}\) and freezing continuous principal coefficients on smaller balls. Interior estimates then give smoothness as in
Lemma~\ref{lem:ricci-radius}. The smooth metric is flat.

In fact \(m_i\to h\) strongly in \(W^{2,2}\) on smaller interiors.
Subtract the harmonic Ricci equations and localize with another
cutoff. The operator on \(m_i-h\) has small principal perturbation
and first-order coefficients bounded by \(C(|Dm_i|+|Dh|)\),
small in \(L^4\). The forcing tends to zero in \(L^2\):
the Ricci norm vanishes, coefficient differences multiplying the
smooth \(D^2h,(Dh)^2\) converge by dominated convergence, and
cutoff commutators with \(D(m_i-h)\) converge strongly in \(L^2\).
For the drift commutator use boundedness in \(L^4\) of the drift
and strong \(L^4\) convergence of \(m_i-h\). Absorption at \(s=2\)
gives the claim, hence strong \(W^{1,4}\).

It remains to place a whole spherical cut in these coordinates.
On \(A(1/4,4)\), integrate an oriented parallel orthonormal coframe
of \(h\) to a developing map \(F\) into \(\mathbb R^4\). The annulus
is simply connected. Its derivative is within \(C\varepsilon_*\)
of \(SO(4)\), since \((DF)^tDF=h\). The rigidity estimate on this
fixed Lipschitz domain gives a rotation \(Q\) with
\(\|DF-Q\|_2\le C\varepsilon_*\)
\cite[Theorem 3.1]{FrieseckeJamesMueller2002Rigidity}.
After translation, Poincar\'e and the uniform Lipschitz bound show
\(F-Qx\) is uniformly small on a slightly smaller closed annulus
as \(\varepsilon_*\downarrow0\). For example a fixed lower bound
on its supremum there forces a fixed positive \(L^2\) norm on an
interior ball. Choose \(\varepsilon_*\) so the degree on
\(A(1/3,3)\) is one at every point of \(A(1/2,2)\), by the
straight homotopy to \(Qx\) on the boundary. Positive Jacobian
gives exactly one preimage there. The pulled-back unit sphere
separates the original inner and outer faces: every crossing of
the trimmed annulus goes from smaller to larger than unit
developed radius. Thus both original charged sets retain their
full masses. This argument uses only a bounded trim; if a wider
fixed log interval is desired, the dilation of the radial splice
profiles below is fixed in advance.

In the developed chart on \(A(1/2,2)\) the metrics converge strongly
in \(W^{1,4}\) to the Euclidean metric and are uniformly elliptic.
Uniform convergence is unnecessary for the cost here. Interpolate
with the Euclidean metric on the two fixed splice strips, leaving
an exactly flat middle strip. For the chiral density matrix the
difference from the cutoff times the incoming matrix has \(L^2\)
norm at most a constant times
\[
 \|\Rm_{m_i}\|_2+\|Dm_i\|_4^2+\|Dm_i\|_2+
                       \|m_i-\delta_{\rm Euc}\|_2=o(1).
\]
Indeed the projection-change term is bounded by the first term
using ellipticity, and the product rule in the all-lowered
curvature formula gives the rest. On each splice separately the same bound gives an individual excess \(o(1)\) above the incoming raw side. The bound holds uniformly along
the cutoff interpolation on the closed original type. Multiplying
conformally to obtain the exact round product costs zero chiral
energy. The path stays below \eqref{eq:zero-wall}. Stretch the
exact product and cut as in Proposition~\ref{prop:short-cut}, choosing
a fixed buffer length for the band margin before the source limit.
Both caps have positive Yamabe sign and their combined excess above
the incoming energy is \(o(1)\). The middle sphere is smoothly
marked and both complete raw sides have positive limiting mass.
This contradicts the fixed-type cut inequality.
\end{proof}

\begin{corollary}[A physical annulus located by one point]\label{cor:anchored-annulus}
Let \(x_i\) attain the smaller physical chiral quantile radius \(a_i=r_i(\theta)\)
on the separated subsequence, with \(0<\theta<1/2\) and median radius
\(b_i=r_i(1/2)\). Fix \(A<\infty\) and \(\Lambda\ge1\).
There is a finite \(L>8\), depending only on \(A,\Lambda\) and the
lower Yamabe bound \(y_0\), for which the following data are impossible
on the fixed nonstandard type at its positive-class infimum.
Suppose \(a_i/\rho_i\to0\), \(\rho_i/b_i\to0\), and there are oriented
smooth embeddings \(\Phi_i:A(1/8,L)\to M\). Write
\(S_i(r)=\Phi_i(\{|z|=r\})\) and \(m_i=\rho_i^{-2}\Phi_i^*g_i\).
Assume
\[
 d_{g_i}(x_i,S_i(1))\le A\rho_i,\qquad
 \Lambda^{-1}\delta_{\rm Euc}\le m_i\le\Lambda\delta_{\rm Euc},\qquad
 \int_{A(1/8,L)}|\Rm_{m_i}|^2\,dV_{m_i}\longrightarrow0 .
\]
On the fixed central subannulus \(A(1/8,8)\) assume the small distortion
and harmonic or differentiated coordinate-tension hypotheses of
Proposition~\ref{prop:critical-annulus}. Then the middle sphere separates
the entire smaller ball from the entire median exterior. No initial
identification of either complementary region is needed.
\end{corollary}

\begin{proof}
Use physical distances and volumes throughout. The Yamabe bound and the
ball cutoff iteration in Corollary~\ref{cor:energy-radius} give
\(\operatorname{Vol} B(z,r)\ge c(y_0)r^4\) for \(0<r\le1\). Thus
the physical diameters are bounded by packing along a minimizing segment,
so \(\rho_i\to0\). Suppress \(i\) temporarily. The connected, two-sided
embedded sphere \(S(1)\) separates since \(H_1(M;\mathbb Z)=0\).
Its two sides are distinguished by the two radial directions in the chart.
Let \(U\) be either side with physical volume at most one half.
On that side define \(f\) to rise from zero at \(S(1)\) to one across
the coordinate strip from radius \(1\) to \(2\) or from \(1\) to
\(1/2\), according to the side. Extend by one on the rest of \(U\)
and by zero on the other side. The radial linear ramp gives
\(\int|df|^2\le C(\Lambda)\rho^2\).
Since the scalar curvature equals the Yamabe constant \(Y\) at volume one,
\[
 (1-2^{-1/2})Y\|f\|_4^2\le6\int|df|^2 .
\]
Hence the plateau has volume at most \(C(y_0,\Lambda)\rho^4\).
Adding the ramp volume gives the same bound for all of \(U\).
A physical ball about \(z\in U\) of radius
\(\min\{1,d(z,S(1))\}\) lies in \(U\). Noncollapse shows that
\(d(z,S(1))\le C'\rho\) for all large \(i\).

If the coordinate outer side had volume at most one half, it would
contain the embedded shell \(2<|z|<L/2\) of physical volume at least
\(\Lambda^{-2}|B^4|((L/2)^4-16)\rho^4\). Choose \(L\) so this
coefficient exceeds the preceding upper bound, and so
\((L/2-1)/\sqrt\Lambda>A\). These choices precede the limit in \(i\).
The inner side is therefore the smaller one. Its every point is within
\((A+C'+\pi\sqrt\Lambda)\rho\) of \(x\), using the intrinsic
diameter bound on the coordinate unit sphere. In particular it misses
the entire median exterior for all large \(i\).

For completeness, the location of the charged center also follows
from the chart tests. Every point of the closed coordinate band
\(1/4\le |z|\le L/2\) has ambient distance at least \(c(\Lambda)\rho\)
from the complement of the chart, by first exit across a buffered
coordinate sphere. If \(B(x,a)\) met this band, the whole ball would
lie in the chart since \(2a=o(\rho)\). Its chiral mass
\(\theta e_i\) is bounded below, contrary to the vanishing full
curvature integral. Thus \(x\) is outside this band. If \(x\) were
on the outer side of \(S(1)\), a minimizing path to \(S(1)\) would
cross the embedded band from \(S(L/2)\) to \(S(1)\). Between its last
hit on \(S(L/2)\) and first hit on \(S(1)\) it stays in that band,
costing at least \((L/2-1)\rho/\sqrt\Lambda\), contrary to the anchor.
The connected inner ball is on the inner side. The trim through radius
four lies within \(C(A,\Lambda)\rho\) of \(x\) by chart paths and thus
misses the median exterior as well. On the two source sides the raw
chiral masses are consequently at least \(\theta e_i\) and \(e_i/2\).
The central splices and positive caps in the proof of
Proposition~\ref{prop:critical-annulus} now contradict the one-cut
bound. All energy estimates here concern the single selected sequence.
\end{proof}

Retain the wide chart, scale, curvature and anchor hypotheses of the corollary. The same marking proof applies if the central metric converges uniformly
and strongly in \(W^{1,4}\) to the Euclidean metric in the supplied
chart. Proposition~\ref{prop:short-cut} then furnishes the splices
directly. In particular a single marked Euclidean annulus on one
component of a physical limit suffices if it has the stated width,
scale, anchor and source convergence. Other components of a distance
shell need no separate parametrizations for this test.
